% ==============================================================================
% Paper 87 (DE): Die Bateman-Horn-Vermutung: Primzahlwerte von Polynomen
% Autor: Michael Fuhrmann
% Projekt: Specialist Mathematics Project, Build 167
% Datum: 2026-03-12
% ==============================================================================

\documentclass[12pt,a4paper]{amsart}

\usepackage{amsmath,amssymb,amsthm,mathtools,enumitem,booktabs,array,hyperref}
\usepackage[utf8]{inputenc}
\usepackage[T1]{fontenc}
\usepackage[ngerman]{babel}
\usepackage{geometry}
\geometry{margin=2.5cm}

% --------------------------------------------------------------------------
% Theorem-Umgebungen (Deutsch)
% --------------------------------------------------------------------------
\newtheorem{theorem}{Theorem}[section]
\newtheorem{lemma}[theorem]{Lemma}
\newtheorem{corollary}[theorem]{Korollar}
\newtheorem{proposition}[theorem]{Proposition}
\newtheorem{satz}[theorem]{Satz}
\newtheorem{korollar}[theorem]{Korollar}
\theoremstyle{definition}
\newtheorem{definition}[theorem]{Definition}
\newtheorem{definition_de}[theorem]{Definition}
\newtheorem{example}[theorem]{Beispiel}
\theoremstyle{remark}
\newtheorem{remark}[theorem]{Bemerkung}
\newtheorem{bemerkung}[theorem]{Bemerkung}
\newtheorem{conjecture}[theorem]{Vermutung}
\newtheorem{vermutung}[theorem]{Vermutung}

% --------------------------------------------------------------------------
% Makros
% --------------------------------------------------------------------------
\newcommand{\Z}{\mathbb{Z}}
\newcommand{\Q}{\mathbb{Q}}
\newcommand{\R}{\mathbb{R}}
\newcommand{\C}{\mathbb{C}}
\newcommand{\F}{\mathbb{F}}
\newcommand{\N}{\mathbb{N}}
\newcommand{\li}{\mathrm{li}}
\DeclareMathOperator{\ord}{ord}
\newcommand{\BH}{\mathfrak{S}}

\hypersetup{
  colorlinks=true,
  linkcolor=blue,
  citecolor=blue,
  urlcolor=blue,
  pdftitle={Die Bateman-Horn-Vermutung: Primzahlwerte von Polynomen},
  pdfauthor={Michael Fuhrmann}
}

% ==============================================================================
\begin{document}
% ==============================================================================

\title{Die Bateman-Horn-Vermutung:\\
  Primzahlwerte von Polynomen}

\author{Michael Fuhrmann}
\address{Specialist Mathematics Project, Build 167}
\email{specialist-maths@project.local}
\date{2026-03-12}

\begin{abstract}
Die Bateman-Horn-Vermutung, 1962 formuliert, liefert eine quantitative Vorhersage für
die Anzahl der ganzen Zahlen $n \leq x$, für die $f_1(n), f_2(n), \ldots, f_k(n)$
gleichzeitig Primzahlen sind, wobei $f_1,\ldots,f_k \in \Z[x]$ paarweise verschiedene
irreduzible Polynome sind, die eine offensichtlich notwendige Bedingung
(\emph{Zulässigkeit}) erfüllen.
Die vermutete Anzahl ist
\[
  \#\{n \leq x : f_i(n) \text{ prim für alle } i\} \;\sim\;
  \BH(f_1,\ldots,f_k) \cdot \frac{x}{(\log x)^k},
\]
wobei $\BH(f_1,\ldots,f_k)$ die \emph{singuläre Reihe} (ein Euler-Produkt über alle
Primzahlen) ist.
Diese Vermutung vereint viele klassische Konvekturen: die Bunyakovsky-Vermutung
($k=1$), die Primzahlzwillings-Vermutung ($f_1 = n, f_2 = n+2$), Sophie-Germain-Primzahlen,
Dicksons Vermutung und Schinzels Hypothese H.
Bewiesene Spezialfälle sind: Dirichlets Satz über Primzahlen in arithmetischen
Folgen (lineare $f_i$), der Green-Tao-Satz und der Maynard-Tao-Satz
(beschränkte Primzahllücken).
Die allgemeine Bateman-Horn-Vermutung ist offen.
\end{abstract}

\maketitle

\tableofcontents

% ==============================================================================
\section{Einleitung}
% ==============================================================================

Eine grundlegende Frage der Zahlentheorie ist: Für welche Polynome $f \in \Z[x]$ gibt
es unendlich viele $n \in \N$ mit $f(n)$ prim? Allgemeiner: Wann sind für ein System
$f_1, \ldots, f_k \in \Z[x]$ unendlich viele $n$ mit $f_1(n),\ldots,f_k(n)$ alle prim,
und wie oft kommt das vor?

Die Bateman-Horn-Vermutung \cite{BH1962} gibt eine einheitliche quantitative Antwort:
Unter einer natürlichen lokalen Bedingung (Zulässigkeit) werden solche $n$ durch ein
Euler-Produkt (die singuläre Reihe) beschrieben.

\subsection{Historischer Hintergrund}

Das Problem von Primzahlwerten bei Polynomen hat eine lange Geschichte:
\begin{itemize}
  \item \textbf{Dirichlet (1837)}: Beweis, dass $an+b$ ($\gcd(a,b)=1$) unendlich viele
    Primzahlen darstellt. Dies ist der einzige vollständig gelöste lineare Fall.
  \item \textbf{Bunyakovsky (1857)}: Vermutung, dass jedes irreduzible $f \in \Z[x]$ mit
    positivem Leitkoeffizient und ohne feste Primteiler unendlich viele Primzahlen
    darstellt. Offen für $\deg f \geq 2$.
  \item \textbf{Hardy-Littlewood (1923)} \cite{HL1923}: Quantitative Vermutung für
    Primzahlpaare und Primzahlkonstellationen.
  \item \textbf{Dickson (1904)} \cite{Dickson1904}: Vermutung über zulässige lineare
    Systeme und Primzahl-Tupel.
  \item \textbf{Schinzel-Sierpiński (1958)} \cite{Schinzel1958}: Hypothese H vereinigt
    alle Polynomial-Vermutungen.
  \item \textbf{Bateman-Horn (1962)} \cite{BH1962}: Quantitative Version von Hypothese H.
\end{itemize}

% ==============================================================================
\section{Grundlagen und Definitionen}\label{sec:grundlagen}
% ==============================================================================

\begin{definition_de}[Zulässiges System]
  Ein System paarweise verschiedener, nicht-konstanter Polynome $f_1, \ldots, f_k \in \Z[x]$
  (mit positivem Leitkoeffizient) heißt \emph{zulässig}, wenn für jede Primzahl $p$
  ein ganzzahliges $n_0$ existiert mit
  $p \nmid f_1(n_0) f_2(n_0) \cdots f_k(n_0)$.
\end{definition_de}

\begin{bemerkung}
  Zulässigkeit ist eine notwendige Bedingung: Gibt es $p$ mit $p | f_i(n)$ für alle
  $n \in \Z$, dann kann $f_i(n)$ nur endlich oft prim sein.
\end{bemerkung}

\begin{example}
  \begin{enumerate}[label=(\roman*)]
    \item $\{n, n+2\}$ ist zulässig: Für $p=2$ ist genau eine der Zahlen $n, n+2$ ungerade;
      für $p\geq 3$ sind $0$ und $-2$ verschiedene Restklassen modulo $p$.
    \item $\{n, n+1\}$ ist \emph{nicht} zulässig: Eine der Zahlen ist stets gerade.
    \item $\{n^2+1\}$ ist zulässig: Für $p \equiv 3 \pmod{4}$ hat $n^2 \equiv -1 \pmod{p}$
      keine Lösung; für $p \equiv 1 \pmod{4}$ nur die Hälfte der Restklassen.
  \end{enumerate}
\end{example}

% ==============================================================================
\section{Die singuläre Reihe}\label{sec:singulaere_reihe}
% ==============================================================================

\begin{definition_de}[Singuläre Reihe]
  Für ein zulässiges System $\{f_1,\ldots,f_k\}$ irreduzibler Polynome über $\Z$ sei
  \[
    \BH(f_1,\ldots,f_k) \;:=\;
    \prod_p \left(1 - \frac{1}{p}\right)^{-k} \left(1 - \frac{\omega(p)}{p}\right),
  \]
  wo $\omega(p)$ die Anzahl der Lösungen von $f_1(n)\cdots f_k(n) \equiv 0 \pmod{p}$
  in $\Z/p\Z$ ist.
\end{definition_de}

\begin{proposition}
  Die singuläre Reihe $\BH(f_1,\ldots,f_k)$ konvergiert absolut zu einer positiven
  reellen Zahl, sofern das System zulässig ist.
\end{proposition}

% ==============================================================================
\section{Die Bateman-Horn-Vermutung}\label{sec:vermutung}
% ==============================================================================

\begin{vermutung}[Bateman-Horn \cite{BH1962}]\label{verm:BH}
  Seien $f_1, \ldots, f_k \in \Z[x]$ paarweise verschiedene irreduzible Polynome
  (mit positiven Leitkoeffizienten $a_1, \ldots, a_k$) und sei das System zulässig.
  Dann gilt:
  \[
    \#\{1 \leq n \leq x : f_1(n),\ldots,f_k(n) \text{ alle prim}\}
    \;\sim\;
    \frac{\BH(f_1,\ldots,f_k)}{\prod_{i=1}^k a_i}
    \cdot \frac{x}{(\log x)^k}
    \quad (x \to \infty).
  \]
\end{vermutung}

\begin{bemerkung}
  Bateman und Horn leiteten diese Vermutung her, indem sie annahmen, dass die Primzahlen
  $f_i(n)$ sich ``unabhängig'' verhalten, außer für die offensichtlichen
  Kongruenz-Obstruktionen bei jeder Primzahl $p$, die in der singulären Reihe
  $\BH$ erfasst werden.
\end{bemerkung}

% ==============================================================================
\section{Spezialfälle}\label{sec:spezialfaelle}
% ==============================================================================

\subsection{Bunyakovsky-Vermutung ($k=1$)}

\begin{vermutung}[Bunyakovsky, 1857 \cite{Bunyakovsky1857}]\label{verm:bunyakovsky}
  Sei $f \in \Z[x]$ irreduzibel mit positivem Leitkoeffizient und $\omega(p) < p$ für
  alle Primzahlen $p$. Dann ist $f(n)$ für unendlich viele $n$ prim.
\end{vermutung}

\begin{example}[$n^2+1$]
  Für $f(n) = n^2+1$ ist die singuläre Reihe das Landau-Ramanujan-ähnliche Produkt
  $\approx 1{,}3203\ldots$. Die Vermutung sagt unendlich viele Primzahlen der Form
  $n^2+1$ voraus -- offen.
\end{example}

\subsection{Primzahlzwillings-Vermutung ($k=2$, $f_1=n$, $f_2=n+2$)}

\begin{vermutung}[Primzahlzwillings-Vermutung]\label{verm:zwillinge}
  Es gibt unendlich viele Primzahlen $p$ mit $p+2$ prim.
\end{vermutung}

\begin{satz}[Hardy-Littlewood-Vermutung B, Spezialfall von Bateman-Horn]
  Unter der Bateman-Horn-Vermutung gilt:
  \[
    \#\{p \leq x : p+2 \text{ prim}\} \;\sim\;
    2C_2 \frac{x}{(\log x)^2},
  \]
  wobei $C_2 = \prod_{p \geq 3} \frac{p(p-2)}{(p-1)^2} \approx 0{,}6601618\ldots$
  die \emph{Primzahlzwillings-Konstante} ist.
\end{satz}

\subsection{Sophie-Germain-Primzahlen ($f_1 = n$, $f_2 = 2n+1$)}

\begin{vermutung}[Unendlichkeit der Sophie-Germain-Primzahlen]\label{verm:sg}
  Es gibt unendlich viele Primzahlen $p$ mit $2p+1$ prim.
\end{vermutung}

Die Bateman-Horn-Vorhersage ergibt den gleichen Ausdruck wie für Primzahlzwillinge
(durch Symmetrie der singulären Reihe).

\subsection{Eulers Primzahl-erzeugendes Polynom}

\begin{example}[$n^2+n+41$]
  Das Polynom $f(n) = n^2+n+41$ ist für $n = 0, 1, 2, \ldots, 39$ stets prim
  (Euler, 1772). Die Diskriminante $\Delta = 1-164 = -163$ hat Klassenzahl $1$,
  wodurch $\omega(p) = 0$ für alle $p \leq 41$. Bateman-Horn sagt unendlich viele
  Primzahlwerte voraus -- offen.
\end{example}

% ==============================================================================
\section{Schinzels Hypothese H}\label{sec:schinzel}
% ==============================================================================

\begin{vermutung}[Schinzels Hypothese H \cite{Schinzel1958}]
  \label{verm:schinzel}
  Seien $f_1, \ldots, f_k \in \Z[x]$ irreduzible Polynome mit positivem Leitkoeffizient,
  die das Zulässigkeitskriterium erfüllen.
  Dann gibt es unendlich viele $n \in \N$, für die alle $f_1(n), \ldots, f_k(n)$ prim sind.
\end{vermutung}

\begin{bemerkung}
  Hypothese H ist die qualitative Version der Bateman-Horn-Vermutung.
  Sie impliziert insbesondere:
  \begin{itemize}
    \item Die Primzahlzwillings-Vermutung (Vermutung~\ref{verm:zwillinge});
    \item Unendlich viele Sophie-Germain-Primzahlen (Vermutung~\ref{verm:sg});
    \item Unendlich viele Primzahlen der Form $n^2+1$;
    \item Dicksons Vermutung für Primzahl-Konstellationen.
  \end{itemize}
\end{bemerkung}

% ==============================================================================
\section{Hardy-Littlewood-Vermutungen}\label{sec:hardy_lw}
% ==============================================================================

\begin{vermutung}[Hardy-Littlewood Vermutung A \cite{HL1923}]
  \label{verm:HLA}
  Die Anzahl der Primzahl-Konstellationen des Typs
  $\{n, n+h_1, \ldots, n+h_{k-1}\}$ (für ein zulässiges Tupel) bis $x$ erfüllt
  \[
    \pi(x; h_1,\ldots,h_{k-1}) \;\sim\; \BH_k \cdot \frac{x}{(\log x)^k}.
  \]
\end{vermutung}

\begin{vermutung}[Hardy-Littlewood Vermutung B (Primzahlpaare)]
  \[
    \#\{n \leq x : n \text{ und } n+2k \text{ beide prim}\}
    \;\sim\; 2C_2 \prod_{p|k, p>2} \frac{p-1}{p-2} \cdot \frac{x}{(\log x)^2}.
  \]
\end{vermutung}

% ==============================================================================
\section{Dicksons Vermutung}\label{sec:dickson}
% ==============================================================================

\begin{vermutung}[Dickson, 1904 \cite{Dickson1904}]
  \label{verm:dickson}
  Seien $a_1 n + b_1, \ldots, a_k n + b_k$ lineare Polynome mit positiven
  Leitkoeffizienten und $b_i \in \Z$, die das Zulässigkeitskriterium erfüllen.
  Dann gibt es unendlich viele $n$, für die alle $a_i n + b_i$ prim sind.
\end{vermutung}

\begin{bemerkung}
  Dicksons Vermutung ist der lineare Spezialfall von Schinzels Hypothese H.
  Spezialfälle:
  \begin{itemize}
    \item $k=1$: Dirichlets Satz (bewiesen, Satz~\ref{satz:dirichlet}).
    \item $k=2$, $\{n, n+2\}$: Primzahlzwillings-Vermutung (offen).
    \item $k=2$, $\{n, 2n+1\}$: Sophie-Germain-Primzahlen (offen).
  \end{itemize}
\end{bemerkung}

% ==============================================================================
\section{Dirichlets Satz: Der bewiesene lineare Fall}\label{sec:dirichlet}
% ==============================================================================

\begin{satz}[Dirichlet, 1837 \cite{Dirichlet1837}]\label{satz:dirichlet}
  Seien $a, b \in \Z$ mit $\gcd(a,b) = 1$ und $a > 0$. Dann gibt es unendlich viele
  Primzahlen der Form $an+b$, $n \geq 1$. Genauer:
  \[
    \#\{p \leq x : p \equiv b \pmod{a}\} \;\sim\; \frac{1}{\phi(a)} \cdot \frac{x}{\log x},
  \]
  wobei $\phi(a)$ die Eulersche Phi-Funktion ist.
\end{satz}

\begin{bemerkung}
  Dirichlets Satz ist der Spezialfall $k=1$, $f(n) = an+b$ der Bateman-Horn-Vermutung.
  Er folgt aus der Nicht-Nullstellen-Aussage $L(1,\chi) \neq 0$ für nicht-hauptcharaktere
  Dirichlet-Charaktere $\chi$ modulo $a$.
\end{bemerkung}

\begin{proof}[Beweisidee]
  Schreibe $\pi(x;a,b) = \frac{1}{\phi(a)} \sum_{\chi \pmod{a}} \bar{\chi}(b) \pi(x,\chi)$,
  wo $\pi(x,\chi) = \sum_{p \leq x} \chi(p)$.
  Für den Hauptcharakter $\chi_0$ gilt $\pi(x,\chi_0) \sim x/\log x$ (PNT).
  Für $\chi \neq \chi_0$ gilt $\pi(x,\chi) = o(x/\log x)$ wegen $L(1,\chi) \neq 0$.
  Summieren ergibt $\pi(x;a,b) \sim x / (\phi(a) \log x)$.
\end{proof}

% ==============================================================================
\section{Sieb-Oberschranken}\label{sec:sieb}
% ==============================================================================

\begin{satz}[Bruns Satz, 1919 \cite{Brun1919}]
  Die Summe der Kehrwerte der Primzahlzwillinge konvergiert:
  \[
    \sum_{p,p+2 \text{ beide prim}} \left(\frac{1}{p} + \frac{1}{p+2}\right) < \infty.
  \]
\end{satz}

\begin{bemerkung}
  Bruns Satz beweist nicht die Unendlichkeit, schließt sie aber auch nicht aus.
  Im Gegensatz dazu divergiert $\sum_p 1/p$ (Euler).
\end{bemerkung}

\begin{satz}[Chen, 1966/1973 \cite{Chen1973}]
  Es gibt unendlich viele Primzahlen $p$ mit $p+2$ entweder prim oder Produkt zweier
  Primzahlen (\emph{Halbprimzahl}).
  Allgemeiner: Jede gerade Zahl $2k$ lässt sich als $p + m$ schreiben, wo $p$ prim und
  $m$ höchstens zwei Primfaktoren hat.
\end{satz}

\begin{satz}[Oberschranke via Selberg-Sieb]
  Für jedes zulässige System $\{f_1,\ldots,f_k\}$ linearer Polynome:
  \[
    \#\{n \leq x : f_i(n) \text{ alle prim}\} \;\leq\;
    C_k \cdot \frac{\BH(f_1,\ldots,f_k) \cdot x}{(\log x)^k}
  \]
  für eine explizite Konstante $C_k > 1$.
  Die Bateman-Horn-Vermutung besagt $C_k = 1 + o(1)$.
\end{satz}

% ==============================================================================
\section{Green-Tao: Primzahlen in arithmetischen Folgen}\label{sec:green_tao}
% ==============================================================================

\begin{satz}[Green-Tao, 2008 \cite{GT2008}]\label{satz:green_tao}
  Die Primzahlen enthalten arithmetische Folgen beliebiger Länge.
  Für jedes $k \geq 1$ gibt es unendlich viele arithmetische Folgen
  $\{a, a+d, a+2d, \ldots, a+(k-1)d\}$ der Länge $k$ (mit $d > 0$), die nur aus
  Primzahlen bestehen.
\end{satz}

\begin{bemerkung}
  Der Green-Tao-Satz beweist die \emph{Unendlichkeit} (konsistent mit Bateman-Horn
  für lineare Tupel), aber nicht die volle asymptotische Formel.
\end{bemerkung}

\begin{bemerkung}[Beweismethoden]
  Der Beweis nutzt:
  \begin{enumerate}[label=(\roman*)]
    \item Den Satz von Szemerédi (arithmetische Folgen in dichten Mengen);
    \item Ein \emph{Transferenzprinzip}: Primzahlen sind „pseudozufällig" genug;
    \item Das verallgemeinerte von-Neumann-Theorem (Green-Tao-Ziegler).
  \end{enumerate}
\end{bemerkung}

\begin{satz}[Quantitativer Green-Tao, Green-Tao-Ziegler \cite{GTZ2012}]
  Die Anzahl der $k$-gliedrigen arithmetischen Primzahlfolgen bis $x$ ist
  $\sim \BH_k \cdot x^2 / (\log x)^k$,
  konsistent mit der Bateman-Horn-Vermutung.
\end{satz}

% ==============================================================================
\section{Maynard-Tao: Beschränkte Primzahllücken}\label{sec:maynard}
% ==============================================================================

\begin{satz}[Zhang, 2013 \cite{Zhang2013}]
  Es gibt unendlich viele Paare von Primzahlen $(p,q)$ mit $|p-q| \leq 70\,000\,000$.
\end{satz}

\begin{satz}[Maynard, 2015 \cite{Maynard2015}; Tao (unabhängig, 2013)]
  \label{satz:maynard}
  Für jedes $k \geq 2$ gilt:
  \[
    \liminf_{n\to\infty}(p_{n+1} - p_n) \;\leq\; 246
  \]
  (Polymath8b \cite{Polymath2014}). Allgemeiner:
  Für jedes $m \geq 1$ gilt $\liminf_{n\to\infty}(p_{n+m} - p_n) < \infty$.
\end{satz}

\begin{bemerkung}
  Der Maynard-Tao-Beweis nutzt ein \emph{mehrdimensionales GPY-Sieb}.
  Er zeigt, dass in einem fest gewählten zulässigen $k$-Tupel unendlich oft
  \emph{mindestens eine} Zahl prim ist -- nicht alle gleichzeitig
  (das würde die Primzahlzwillings-Vermutung liefern).
\end{bemerkung}

\begin{bemerkung}[Abstand zu Bateman-Horn]
  \begin{itemize}
    \item Maynard-Tao: \emph{Mindestens eine} Primzahl in einem zulässigen $k$-Tupel,
      unendlich oft.
    \item Bateman-Horn: \emph{Alle} gleichzeitig prim, unendlich oft und quantitativ.
  \end{itemize}
\end{bemerkung}

% ==============================================================================
\section{Zulässigkeit und Obstruktionen}\label{sec:zulässigkeit}
% ==============================================================================

\begin{proposition}[Zulässigkeitstest für lineare Polynome]
  Das System $\{a_1 n + b_1, \ldots, a_k n + b_k\}$ ist zulässig genau dann, wenn
  für jede Primzahl $p$ die Vereinigung $\{b_i/a_i \pmod{p} : i=1,\ldots,k\}$ nicht
  alle Restklassen von $\Z/p\Z$ abdeckt.
\end{proposition}

\begin{example}[Kleinste zulässige $k$-Tupel]
  \begin{center}
    \begin{tabular}{ll}
      \toprule
      $k$ & Kleinstes zulässiges $k$-Tupel \\
      \midrule
      2 & $(0, 2)$ \\
      3 & $(0, 2, 6)$ \\
      4 & $(0, 2, 6, 8)$ \\
      5 & $(0, 2, 6, 8, 12)$ \\
      6 & $(0, 4, 6, 10, 12, 16)$ \\
      \bottomrule
    \end{tabular}
  \end{center}
\end{example}

% ==============================================================================
\section{Offene Probleme}\label{sec:offene_probleme}
% ==============================================================================

\begin{vermutung}[Bateman-Horn, allgemein; Vermutung~\ref{verm:BH}]
  Offen für alle $k \geq 1$ mit $\deg f_i \geq 2$ und für $k \geq 2$ mit $\deg f_i = 1$.
  Der einzige bewiesene Fall: $k=1$, $\deg f = 1$ (Dirichlet).
\end{vermutung}

\begin{vermutung}[Primzahlzwillings-Vermutung; Vermutung~\ref{verm:zwillinge}]
  Offen. Bestes Ergebnis: Unendlich viele Primzahllücken $\leq 246$
  (Polymath8b \cite{Polymath2014}).
\end{vermutung}

\begin{vermutung}[Unendlichkeit der Primzahlen der Form $n^2+1$]
  Offen. Iwaniec \cite{Iwaniec1978} bewies $\liminf \Omega(n^2+1) \leq 2$
  ($\Omega$ = Anzahl der Primfaktoren mit Vielfachheit), aber kein Primzahlbeweis.
\end{vermutung}

\begin{vermutung}[Bunyakovsky / Hypothese H für $k=1$, $\deg \geq 2$]
  Offen. Kein einziges irreduzibles Polynom vom Grad $\geq 2$ ist bekannt, das unendlich
  viele Primzahlen darstellt.
\end{vermutung}

% ==============================================================================
\section{Fazit}\label{sec:fazit}
% ==============================================================================

Die Bateman-Horn-Vermutung bietet einen eleganten, einheitlichen quantitativen Rahmen
für die Verteilung von Primzahlwerten von Polynom-Systemen.
Ihre Spezialfälle -- Primzahlzwillinge, Sophie-Germain-Primzahlen, Bunyakovsky,
Dickson und Schinzels Hypothese H -- gehören zu den verlockendsten offenen Problemen
der Zahlentheorie.
Der einzige vollständig bewiesene Fall der Bateman-Horn-Familie ist Dirichlets Satz
(lineare Polynome, $k=1$). Der Green-Tao-Satz beweist unendliche arithmetische
Folgen von Primzahlen, und Maynard-Tao etabliert beschränkte Primzahllücken, aber
keines der beiden Ergebnisse löst die allgemeine Vermutung.

Der Abstand zwischen dem Bekannten (Sieb-Oberschranken, Dirichlet, Green-Tao,
Maynard-Tao) und der vollen Bateman-Horn-Vermutung gehört zu den tiefsten
Herausforderungen der analytischen Zahlentheorie.

\begin{thebibliography}{99}

\bibitem{BH1962}
P.\,T.\ Bateman, R.\,A.\ Horn,
\textit{A heuristic asymptotic formula concerning the distribution of prime numbers},
Math.\ Comp.\ \textbf{16} (1962), 363--367.

\bibitem{Brun1919}
V.\ Brun,
\textit{La série $1/5+1/7+1/11+1/13+\cdots$ est convergente},
Bull.\ Sci.\ Math.\ \textbf{43} (1919), 100--104.

\bibitem{Bunyakovsky1857}
V.\ Bunyakovsky,
\textit{Nouveaux théorèmes relatifs à la distinction des nombres premiers},
Mém.\ Acad.\ Sci.\ St.\ Pétersbourg \textbf{6} (1857), 305--329.

\bibitem{Chen1973}
J.\,R.\ Chen,
\textit{On the representation of a larger even integer as the sum of a prime and the
product of at most two primes},
Sci.\ Sinica \textbf{16} (1973), 157--176.

\bibitem{Dickson1904}
L.\,E.\ Dickson,
\textit{A new extension of Dirichlet's theorem on prime numbers},
Messenger Math.\ \textbf{33} (1904), 155--161.

\bibitem{Dirichlet1837}
P.\,G.\,L.\ Dirichlet,
\textit{Beweis des Satzes, daß jede unbegrenzte arithmetische Progression unendlich viele
Primzahlen enthält},
Abh.\ Königl.\ Preuss.\ Akad.\ Wiss.\ (1837), 45--81.

\bibitem{GT2008}
B.\ Green, T.\ Tao,
\textit{The primes contain arbitrarily long arithmetic progressions},
Ann.\ of Math.\ (2) \textbf{167} (2008), Nr.\ 2, 481--547.

\bibitem{GTZ2012}
B.\ Green, T.\ Tao, T.\ Ziegler,
\textit{An inverse theorem for the Gowers $U^{s+1}[N]$-norm},
Ann.\ of Math.\ (2) \textbf{176} (2012), Nr.\ 2, 1231--1372.

\bibitem{HL1923}
G.\,H.\ Hardy, J.\,E.\ Littlewood,
\textit{Some problems of ``Partitio Numerorum'' III},
Acta Math.\ \textbf{44} (1923), 1--70.

\bibitem{Iwaniec1978}
H.\ Iwaniec,
\textit{Almost-primes represented by quadratic polynomials},
Invent.\ Math.\ \textbf{47} (1978), Nr.\ 2, 171--188.

\bibitem{Maynard2015}
J.\ Maynard,
\textit{Small gaps between primes},
Ann.\ of Math.\ (2) \textbf{181} (2015), Nr.\ 1, 383--413.

\bibitem{Polymath2014}
D.\,H.\,J.\ Polymath,
\textit{Variants of the Selberg sieve, and bounded intervals containing many primes},
Res.\ Math.\ Sci.\ \textbf{1} (2014), Art.\ 12.

\bibitem{Schinzel1958}
A.\ Schinzel, W.\ Sierpiński,
\textit{Sur certaines hypothèses concernant les nombres premiers},
Acta Arith.\ \textbf{4} (1958), 185--208.

\bibitem{Zhang2013}
Y.\ Zhang,
\textit{Bounded gaps between primes},
Ann.\ of Math.\ (2) \textbf{179} (2014), Nr.\ 3, 1121--1174.

\end{thebibliography}

\end{document}
